% ISOLATED FUTURE INTEGRATION CANDIDATE. See GATE.md; not a live proof replacement.

\begin{theorem}
\label{thm:ch29-real-class-curve-model-envelope}
Let $K$ be a complete nontrivially valued non-Archimedean field, and let
$C$ be a smooth proper geometrically connected $K$-curve.  Write
$\mathcal D_{\mathbf Q}(C)$ for the rational model functions, namely the
functions induced by formal metrics on positive tensor powers of the
trivial line bundle, divided by that tensor power.  Let
$\theta\in Z^{1,1}(C)$ be a real closed model $(1,1)$-form in the sense of
\cite[\S\S2.3--2.7]{GJKM19}, and put $d=\deg\{\theta\}$.
For $f\in C^0(C^{\mathrm{an}},\mathbf R)$ define
\begin{equation}
\label{eq:ch29-real-rational-model-envelope}
 P^{\mathcal D_{\mathbf Q}}_\theta(f)
 :=\sup\{\varphi\in\mathcal D_{\mathbf Q}(C):
             \theta+\mathrm d\mathrm d^c\varphi\text{ is nef},\quad
             \varphi\leq f\},
 \qquad \sup\varnothing=-\infty.
\end{equation}
Here nefness is numerical nefness on a determining proper model, not
merely nefness of the generic class.  Then the following hold.
\begin{enumerate}
\item If $d>0$, the envelope in
\eqref{eq:ch29-real-rational-model-envelope} is a finite continuous function
for every continuous $f$.  It is a uniform limit of $\theta$-psh rational
model functions lying below $f$.
\item Suppose $d=0$ and an actual function
$\psi\in\mathcal D_{\mathbf Q}(C)$ with
$\theta+\mathrm d\mathrm d^c\psi$ nef exists.  Then
$\theta+\mathrm d\mathrm d^c\psi=0$ in $Z^{1,1}(C)$, any two such
functions differ by a real constant, and for every continuous $f$,
\begin{equation}
\label{eq:ch29-real-degree-zero-envelope}
 P^{\mathcal D_{\mathbf Q}}_\theta(f)
       =\psi+\min_{C^{\mathrm{an}}}(f-\psi).
\end{equation}
This is again a uniform limit of $\theta$-psh rational model functions
lying below $f$.  If no such $\psi$ exists, all these envelopes equal
$-\infty$.
\end{enumerate}
The hypothesis $d>0$ is equivalent to ampleness of the real numerical
class on the curve.  No assertion here identifies
$P^{\mathcal D_{\mathbf Q}}_\theta$ with an envelope over the real span
of model functions, or with an envelope over all singular psh functions.
\end{theorem}

\begin{proof}
Throughout this proof the word model means rational model unless explicitly
qualified.  We use the model-function and real-model-form conventions of
\cite[\S\S2.2--2.8]{GJKM19}; in particular the coefficients of $\theta$
need not be rational and $\theta$ need not come from one line bundle.
A smooth proper curve over a field is projective, so the proper and
projective model results cited below apply to $C$.

\emph{Elementary envelope operations.}
Put $\Gamma_K=-\log|K^\times|$ and
$\Lambda_K=\mathbf Q\Gamma_K$.  Constants in $\Lambda_K$ are model
functions with zero curvature.  This subgroup is dense in $\mathbf R$:
it contains $a\mathbf Q$ for a nonzero $a\in\Gamma_K$.
Model $\theta$-psh functions are closed under finite maxima
\cite[\S2.5]{GJKM19}.  Thus, if one such function exists, translating it
down by a sufficiently large element of $\Lambda_K$ gives a competitor
under every continuous obstacle.  Every resulting envelope is finite,
bounded above by the obstacle and below by a continuous competitor.
Monotonicity and translations by model constants imply, by taking limits
of constants in $\Lambda_K$, that
\begin{equation}
\label{eq:curve-envelope-contraction}
 P^{\mathcal D_{\mathbf Q}}_\theta(f+c)
       =P^{\mathcal D_{\mathbf Q}}_\theta(f)+c,
 \qquad
 \|P^{\mathcal D_{\mathbf Q}}_\theta(f)
       -P^{\mathcal D_{\mathbf Q}}_\theta(g)\|_\infty
       \leq\|f-g\|_\infty.
\end{equation}
Here $c$ is any real number; this identity does not assert that adding
$c$ preserves the model-function class.  For a model function $b$, the
actual bijection $\varphi\mapsto\varphi-b$ of competitor sets gives
\begin{equation}
\label{eq:ch29-real-model-obstacle-translation}
 P^{\mathcal D_{\mathbf Q}}_\theta(b)-b
       =P^{\mathcal D_{\mathbf Q}}_{\theta+\mathrm d\mathrm d^cb}(0).
\end{equation}
These are also \cite[Proposition~2.9(3)--(5)]{GJKM19}.

\emph{The graph calculation over an algebraically closed field.}
For now assume that $K$ is algebraically closed.  Its value group
$\Gamma_K$ is divisible.  Semistable reduction and domination give a
proper strictly semistable model determining any prescribed finite
collection of model functions and real model forms.  For this curve-model
input and its use for an arbitrary real form, see
\cite[proof of Proposition~3.7]{GJKM19} and the semistable domination
result \cite[\S7]{BL85} used there.  A real form is determined by finitely
many model numerical classes, so common domination applies to it too.
Write $G$ for the finite connected skeleton, $V$ for its nonempty vertex
set, and $\tau:C^{\mathrm{an}}\to G$ for its retraction.  Subdivisions
at $\Gamma_K$-rational points give dominating models.  The edge lengths
$\ell_e$ belong to $\Gamma_K$.

Let $d_v\in\mathbf R$ be the degree of $\theta$ on the component
corresponding to $v$.  The special fiber is reduced and the residue field
is algebraically closed, so
\begin{equation}
\label{eq:ch29-real-graph-curvature}
 D=\sum_{v\in V}d_v\delta_v,
 \qquad \sum_{v\in V}d_v=d,
 \qquad
 \theta+\mathrm d\mathrm d^c(F\circ\tau)\text{ is nef}
       \Longleftrightarrow D+\Delta F\geq0
\end{equation}
whenever $F\circ\tau$ is a model function.  The Laplacian has the
outgoing-slope convention; in an edge interior its coefficient at a break
is the right slope minus the left slope.  Formula
\eqref{eq:ch29-real-graph-curvature} is
\cite[Proposition~3.7]{GJKM19}.  The degree identity is the degree
formula on a proper model, applied real-linearly.  In particular, the
condition is convexity on each open edge and
$d_v+\sum_{\nu\in T_vG}\partial_\nu F\geq0$ at each vertex.
We also use precisely the following model dictionary:
if $F$ has rational slopes, all breakpoints have coordinates in
$\Gamma_K$, and all its vertex values are in $\Gamma_K$, then
$F\circ\tau$ is a model function.  Indeed a common positive integer
multiple is integral $\Gamma_K$-affine, and
\cite[Proposition~3.5(3)]{GJKM19} applies.  This use of divisibility is
why the construction is first made over an algebraically closed field.

Assume $d>0$.  Write $N=|V|$.  There is a real function $h$, affine on
each edge, with
\begin{equation}
\label{eq:ch29-real-positive-vertex-margin}
 D+\Delta h=\frac dN\sum_{v\in V}\delta_v.
\end{equation}
To see this, the Laplacian on edge-affine functions is a symmetric
matrix whose quadratic form, with sign reversed, is
$\sum_{e=(v,w)}(x_v-x_w)^2/\ell_e$.  Its kernel consists exactly of
constants, since $G$ is connected.  Its image is therefore the vectors
with coordinate sum zero.  The right side of
\eqref{eq:ch29-real-positive-vertex-margin} minus $D$ has that sum.
Normalize the solution by one vertex value and then shift it by a real
constant so that $h\leq0$.  This is an auxiliary real graph function;
we have not asserted that it is a model function.

We next prove the approximation needed below.  Let $H\leq0$ be real
and affine on every edge, with
$d_v+\sum_\nu\partial_\nu H\geq\eta>0$ at all vertices.  For every
$\varepsilon>0$ there is a model graph function $A\leq0$ with
$D+\Delta A\geq0$ and $\|A-H\|_\infty<\varepsilon$.
If the graph has no edges, approximate its single value from below by
an element of $\Gamma_K$; this proves the assertion in that case.
Otherwise choose a small $\delta>0$ so that the maximum valence times
$\delta$ is less than $\eta/2$, and so that
$\delta\max_e\ell_e$ is as small as desired.  Choose values
$p_v\in\Gamma_K$ sufficiently close to $H(v)$ that every chord slope
\[
 \sigma_e=\frac{p_w-p_v}{\ell_e}
\]
differs from the corresponding slope of $H$ by less than $\delta/3$.
This is possible because there are finitely many positive edge lengths.
For an oriented edge $e=[v,w]$, choose rational numbers
$\alpha_e<\sigma_e<\beta_e$ whose distances from $\sigma_e$ are less
than $\delta/3$, and put
\begin{equation}
\label{eq:ch29-real-rational-breakpoint}
 t_e=\frac{\beta_e\ell_e-(p_w-p_v)}{\beta_e-\alpha_e},
 \qquad 0<t_e<\ell_e.
\end{equation}
Divisibility gives $t_e\in\Gamma_K$.  On the edge, start at $p_v$,
use slope $\alpha_e$ up to $t_e$, and slope $\beta_e$ thereafter.
Equation \eqref{eq:ch29-real-rational-breakpoint} gives endpoint $p_w$.
These functions glue to a continuous function $A_0$ with rational slopes
and values in $\Gamma_K$ at every old and new vertex.  Its interior
Laplacian is nonnegative because $\beta_e>\alpha_e$.  Each outgoing
slope at an old vertex differs from the corresponding slope of $H$ by
less than $\delta$, so the curvature at that vertex is at least
$\eta/2$.  On every edge the difference from $H$ is bounded by the
initial endpoint error plus $\ell_e\delta$.  Thus $A_0$ is uniformly
as close to $H$ as desired and is a model function by the dictionary.
Finally subtract a positive constant $\lambda\in\Gamma_K$ slightly
larger than $\|A_0-H\|_\infty$.  The resulting $A$ is nonpositive,
retains its curvature, and remains arbitrarily close to $H$.

More generally, if $H_0\leq0$ is edge-affine and satisfies
$D+\Delta H_0\geq0$, replace it by
$(1-s)H_0+sh$ with $0<s<1$.  This has vertex margin $sd/N$ by
\eqref{eq:ch29-real-positive-vertex-margin}, is nonpositive, and converges
uniformly to $H_0$ as $s\downarrow0$.  The preceding construction proves
\begin{equation}
\label{eq:ch29-real-real-chord-approximation}
 \text{every such }H_0\circ\tau\text{ is a uniform limit of
 model }\theta\text{-psh functions }A\circ\tau\leq0.
\end{equation}
Applying the construction directly to $h$ also supplies an actual model
competitor.  Thus nonemptiness has been proved, not inferred from generic
nefness or from an envelope theorem.

\emph{Continuity for a model obstacle over the algebraically closed field.}
Choose the semistable model to determine both $\theta$ and the obstacle
$b\in\mathcal D_{\mathbf Q}(C)$, and use
\eqref{eq:ch29-real-model-obstacle-translation}.  Thus the obstacle is zero
and the new form still has degree $d>0$.
Let $\mathcal A$ consist of the nonpositive graph functions $F$ such
that $F\circ\tau$ is a $\theta$-psh model function.
The preceding construction shows $\mathcal A\ne\varnothing$.
For every model $\theta$-psh function $\varphi$,
\cite[Proposition~3.8]{GJKM19} says that $\varphi\circ\tau$ is again
model and $\theta$-psh, and $\varphi\leq\varphi\circ\tau$.
For an obstacle zero this remains a competitor.  Hence
\begin{equation}
\label{eq:curve-envelope-retraction}
 P^{\mathcal D_{\mathbf Q}}_\theta(0)=F_*\circ\tau,
 \qquad F_*:=\sup_{F\in\mathcal A}F.
\end{equation}
The cited retraction proposition is applied to actual competitors and
does not provide their existence.

For $F\in\mathcal A$, let $L(F)$ be the edge-affine interpolation of
its values on $V$.  Edgewise convexity gives $F\leq L(F)\leq0$.
Chording increases the outgoing slope at each endpoint, so
$D+\Delta L(F)\geq0$.  Formula
\eqref{eq:ch29-real-real-chord-approximation} shows $L(F)\leq F_*$;
the chord itself need not be a model function.
Since \(F\leq L(F)\leq F_*\) for every \(F\in\mathcal A\), we have
\begin{equation}
\label{eq:curve-chord-supremum}
                 F_*=\sup_{F\in\mathcal A}L(F).
\end{equation}
This equality does not declare the real chords to be model functions.
The class $\mathcal A$ is closed under finite maxima.  Given
$\varepsilon>0$, choose $F_v\in\mathcal A$ with
$F_v(v)>F_*(v)-\varepsilon$ at each of the finitely many vertices.
For $B=\max_{v\in V}F_v$ we have
\[
 L(F_*)-\varepsilon<L(B)\leq F_*\leq L(F_*).
\]
Here the left inequality follows by affine interpolation of the vertex
inequalities.  For the right inequality, every $F\in\mathcal A$
is convex on each edge and is bounded at its endpoints by $F_*$, so
it lies below the interpolation $L(F_*)$; take the supremum.
Letting $\varepsilon\downarrow0$ gives $F_*=L(F_*)$.
The envelope is consequently continuous, indeed edge-affine after
subtracting the model obstacle.

Model functions are uniformly dense in $C^0(C^{\mathrm{an}})$ by
\cite[Proposition~2.15]{GM19}.  Together with
\eqref{eq:curve-envelope-contraction}, the result for model
obstacles proves continuity for every continuous obstacle over the
algebraically closed field.

\emph{Field descent, including real model forms.}
For a general $K$, put $\widehat{\overline K}=C_K$ and let
$q:C_{C_K}^{\mathrm{an}}\to C^{\mathrm{an}}$ be the projection.
We claim, with empty suprema permitted, that
\begin{equation}
\label{eq:ch29-real-completed-algebraic-envelope-base-change}
 q^*P^{\mathcal D_{\mathbf Q}}_\theta(f)
  =P^{\mathcal D_{\mathbf Q}}_{q^*\theta}(q^*f).
\end{equation}
Pulling back a competitor proves the inequality from left to right.
For the reverse inequality, take an actual model competitor $a$ on
$C_{C_K}$.  Applying \cite[Proposition~2.18(b)]{GM19} to its metric on
the trivial line bundle gives a finite extension $K'/K$ inside
$\overline K$ and a model function $a'$ on $C_{K'}$ pulling back to
$a$.

The form $\theta_{K'}+\mathrm d\mathrm d^ca'$ is nef because its
pullback to $C_K$ is nef.  This implication is numerical and applies to
\emph{real} model classes: put the finitely many classes representing
$\theta_{K'}$ and the model of $a'$ on a common proper model.
Nefness of a model numerical class is unchanged under scalar extension
and under passing to a dominating model.  For scalar extension one tests
on the components of special-fiber curves; their degrees after extension
are positive factors of the original degrees.  For domination,
exceptional curves have pullback degree zero and the projection formula
on a component mapping onto a curve multiplies its degree by the
positive degree of that map.  These identities extend from line bundles
to real numerical classes by linearity.  Thus nonnegative degrees after
pullback imply nonnegative degrees before pullback.  This is the real-form
base-change invariance used explicitly in
\cite[proof of Proposition~3.7]{GJKM19}; it is not an inference from a
statement only about one metrized line bundle.  The inequality
$a'\leq f_{K'}$ follows from surjectivity of the analytic base-change
map.

Choose a finite normal extension $K''/K$ containing $K'$.  Pulling $a'$
to $K''$ and applying \cite[Proposition~2.11]{GJKM19} gives
\[
 a_{K''}\leq P^{\mathcal D_{\mathbf Q}}_{\theta_{K''}}(f_{K''})
           =q_{K''}^*P^{\mathcal D_{\mathbf Q}}_\theta(f).
\]
This proposition includes inseparable normal extensions; a Galois-only
quotient argument would not suffice.  Pull back once more to $C_K$ and
take the supremum over $a$ to obtain
\eqref{eq:ch29-real-completed-algebraic-envelope-base-change}.
Geometric connectedness ensures the curve upstairs remains connected.
If $d>0$, its envelope is finite and continuous by the graph proof.
The map $q$ is continuous and surjective from a compact space to a
Hausdorff space and hence is a quotient map.  Equation
\eqref{eq:ch29-real-completed-algebraic-envelope-base-change} therefore
proves finiteness and continuity downstairs as well.  In particular the
downstairs competitor family is nonempty.

\emph{The degree-zero case.}
Suppose now $d=0$ and $\psi$ is an actual model $\theta$-psh function.
Over $C_K$, choose a common strictly semistable model determining
$q^*\theta$, $q^*\psi$, and any other model competitor $q^*\varphi$.
All component degrees of
$q^*\theta+\mathrm d\mathrm d^c(q^*\psi)$ are nonnegative and their
sum is zero.  Each is therefore zero; the same holds for $\varphi$.
On a proper curve model the vertical integral curves are exactly the
irreducible special-fiber components.  Vanishing of these degrees says
that both real relative numerical classes are zero.
The reflection argument above gives
$\theta+\mathrm d\mathrm d^c\psi=0$ downstairs too.
On the common connected skeleton the function
$q^*(\varphi-\psi)=U\circ\tau$ has $\Delta U=0$.
After subdividing at its finitely many breaks, graph integration by parts
gives
\[
 0=-\int_G U\,\mathrm d\Delta U
   =\sum_e\int_e|U'|^2\,\mathrm ds.
\]
Thus $U$ is constant on every edge and, by connectedness, on $G$.
Retraction and surjectivity of $q$ imply that $\varphi-\psi$ is a
constant on $C^{\mathrm{an}}$.

Put $c_*:=\min_{C^{\mathrm{an}}}(f-\psi)$; the minimum exists by
compactness.  Any competitor below $f$ equals $\psi+c$ for a constant
$c\leq c_*$, proving the upper bound in
\eqref{eq:ch29-real-degree-zero-envelope}.
Choose $c_j\in\Lambda_K$ with $c_j\leq c_*$ and $c_j\to c_*$.
Then $\psi+c_j$ are actual model $\theta$-psh competitors below $f$,
and converge uniformly to $\psi+c_*$.  This proves the reverse bound
and the required uniform approximation.  No arbitrary real normalization
of $\psi$ has been declared a model function.  If there is no $\psi$,
the empty-family conclusion follows directly from the definition.

\emph{Uniform approximation in the positive-degree case.}
Finally let $P=P^{\mathcal D_{\mathbf Q}}_\theta(f)$ be any of the
finite continuous envelopes just proved.  At each point choose a model
competitor whose value exceeds $P-\varepsilon$ there.  By continuity this
inequality holds on a neighborhood.  A finite subcover of the compact
curve and the finite maximum of its competitors give a model function
$\varphi_\varepsilon$ with
$P-\varepsilon<\varphi_\varepsilon\leq P\leq f$ everywhere.
Taking $\varepsilon\downarrow0$ proves the claimed uniform
approximation.  This is the compact maximum argument of
\cite[Proposition~2.9(8)]{GJKM19} and does not use an envelope theorem
as an input.
\end{proof}

\paragraph{Curvature conventions.}
We use \(\mathrm d\mathrm d^c\) for Thuillier's operator and
\(\mathrm d'\mathrm d''\) for the CLD operator, with the outgoing-slope
sign.  Curvature is the operator applied to the additive weight
\(-\log\lVert e\rVert\).  A type-\(2\) or type-\(3\) point contributes
a unit Dirac \(\delta_q\); a closed algebraic point \(z\) contributes
\([\kappa(z):K]\delta_z\) to its divisor measure.
The original-field exactness, regularity, and divisor formulas below
are \cite[Proposition~3.3.11, Theorem~3.3.13(ii), and
Proposition~3.3.15]{Thu05}.  Thuillier's standing hypotheses here are a
complete nontrivially valued ground field and a smooth strict analytic
curve, not an algebraically closed field
\cite[opening of Chapter~2, p.~15, and Chapter~3, p.~61]{Thu05}.

\begin{proposition}
\label{prop:ch29-curve-point-reference}
Let \(K\) be a complete nontrivially valued non-Archimedean field, and
let \(C\) be a smooth proper geometrically connected algebraic
\(K\)-curve.  Let \(L\) be a line bundle, \(d=\deg L\), and let
\(q\in C^{\mathrm{an}}\) be of type \(2\) or \(3\).
There is a metric \(\lVert\cdot\rVert_q\) on \(L\), lisse in
Thuillier's sense, with curvature
\begin{equation}
\label{eq:ch29-curve-point-reference-curvature}
            c_1^{\mathrm{Thu}}(L,\lVert\cdot\rVert_q)=d\delta_q.
\end{equation}
It is unique up to multiplication by a positive real constant.  If
\(d\geq0\), its local additive weights are Thuillier-subharmonic.
There is no assertion that this metric is rational formal.
\end{proposition}


\begin{proof}
Choose a nonzero rational section \(s\) of \(L\) and put
\(D=\operatorname{div}(s)\).  The finite signed measure
\(d\delta_q-\delta_D\) has total mass zero.  By
\cite[Theorem~3.3.13(ii)]{Thu05}, there is a Thuillier current
\(T\in D^0(C^{\mathrm{an}})\) with
\begin{equation}
\label{eq:ch29-curve-reference-current}
              \mathrm d\mathrm d^cT=d\delta_q-\delta_D.
\end{equation}
Let \(e\) be an invertible local frame, and write \(s=f_e e\).
The Poincar\'e--Lelong formula gives
\[
 \mathrm d\mathrm d^c\bigl(T+\log|f_e|\bigr)=d\delta_q
\]
on that frame's domain.  The right side is a lisse \(1\)-form in
Thuillier's terminology because \(q\) is of type \(2\) or \(3\).
Current regularity \cite[Proposition~3.3.11]{Thu05} therefore makes
\(T+\log|f_e|\) the current of a unique lisse real function \(h_e\).

On an overlap with \(e'=a e\), one has \(f_{e'}=f_e/a\), hence
\(h_{e'}=h_e-\log|a|\).  This is exactly the metric compatibility
for \(h_e=-\log\lVert e\rVert_q\).  The functions thus define a
lisse metric on the original \(K\)-line bundle, with curvature
\eqref{eq:ch29-curve-point-reference-curvature}.  If \(d\geq0\),
\cite[Theorem~3.4.12(ii)]{Thu05} makes every local weight subharmonic.

For two metrics of the asserted curvature, their additive difference
is a continuous global function with zero \(\mathrm d\mathrm d^c\).
It is harmonic by \cite[Lemma~3.3.12]{Thu05}, and constant by the
maximum principle on the proper connected curve
\cite[Proposition~3.1.1]{Thu05}.  This proves uniqueness.  In particular
the construction is not a descent of an arbitrarily chosen Dirac mass
upstairs, and no residue degree is to be inserted at \(q\).
\end{proof}


\begin{theorem}
\label{thm:curves-arbitrary-field-continuity}
Let \(K\) be a complete nontrivially valued non-Archimedean field,
let \(C\) be a smooth projective geometrically connected \(K\)-curve,
and let \(L\) have degree \(d>0\).  The envelope over semipositive
rational model metrics of every continuous metric on \(L^{\mathrm{an}}\)
is a continuous semipositive metric.

For the equation assertion, suppose a proper strictly semistable formal model
\(\mathfrak C\) of \(C^{\mathrm{an}}\) over \(K^\circ\) is given.
Let \(G\) be its skeleton, \(\tau:C^{\mathrm{an}}\to G\) its
retraction, and \(\mu\) a positive Radon measure supported on \(G\)
with mass \(d\).  Assign the edge corresponding to the node \(p_e\)
the conductance
\begin{equation}
\label{eq:ch29-curve-node-conductance}
                  c_e=[\kappa(p_e):\widetilde K].
\end{equation}
Let \(U_\nu\) denote the normalized weighted potential in
\cref{lem:ch31-green-kernel}.  Choose a vertex \(q\) and the reference metric of
\cref{prop:ch29-curve-point-reference}.  Then
\begin{equation}
\label{eq:ch29-weighted-curve-solution}
 U=U_{\mu-d\delta_q},
 \qquad
 \lVert\cdot\rVert_\mu
       =\lVert\cdot\rVert_q\mathrm e^{-U\circ\tau}
\end{equation}
is a continuous metric with Thuillier curvature \(\mu\), unique up to
positive constant scale among continuous metrics of that curvature.
It is a uniform limit of semipositive rational formal metrics over
\(K\), and its Chambert--Loir and CLD curvature measures both equal
\(\mu\).  The same approximation holds after every complete scalar
extension.

If instead a rational formal reference metric \(\lVert\cdot\rVert_0\)
is already determined on \(\mathfrak C\), with curvature
\(\mu_0=\sum_v d_v\delta_v\), the solution can equally be written
\begin{equation}
\label{eq:ch29-curve-model-reference-solution}
 \lVert\cdot\rVert_0\mathrm e^{-U_{\mu-\mu_0}\circ\tau}.
\end{equation}
Here \(d_v=m^{-1}\deg_{\widetilde K}(\mathfrak L|_{C_v})\) if
\(\mathfrak L\) extends \(L^{\otimes m}\); the degree is over
\(\widetilde K\), not over the component's constant field.
No model extension of \(L\) on the originally specified
\(\mathfrak C\) is assumed in \eqref{eq:ch29-weighted-curve-solution}.
For nonatomic \(\mu\) only continuity, not lisse regularity, is asserted.
\end{theorem}

\begin{proof}
For the envelope assertion, choose any rational formal reference
metric on \(L\); one exists by \cite[\S2.1 and Proposition~2.10]{GM19}.
Write the given continuous metric as
\(\lVert\cdot\rVert_b\mathrm e^{-f}\) and put
\(\theta=c_1(L,\lVert\cdot\rVert_b)\).
Apply \cref{thm:ch29-real-class-curve-model-envelope} in degree \(d>0\).
Its rational model competitors are exactly the semipositive rational
model metrics relative to this reference, and its uniform approximation
proves the first assertion.  The reference need not be semipositive.

We now prove the equation assertion independently of envelope continuity.
The weighted graph lemma \cref{lem:ch31-green-kernel}, proved below by
finite-dimensional linear algebra and energy estimates, will be used.

\emph{1. Normalization on the original field.}
The skeleton is a finite connected weighted metric graph; connectedness
is \cite[Remark~2.2.19(3)]{Thu05}.  The weights are exactly
\eqref{eq:ch29-curve-node-conductance}, by
\cite[\S2.2.3, p.~41]{Thu05}.  They are constant along the edge and
the same at its two ends, because both incidences correspond to the
same node.  A subdivision retains this weight.  For every real
piecewise affine \(F\) on \(G\), Thuillier's exact comparison is
\begin{equation}
\label{eq:ch29-curve-retraction-current}
  \mathrm d\mathrm d^c(F\circ\tau)=\iota_*\Delta_c F,
\end{equation}
where \(\iota:G\hookrightarrow C^{\mathrm{an}}\).
Here \(\Delta_c\) denotes the weighted operator of
\eqref{eq:ch31-graph-laplacian}, with \(G\) in place of \(\Gamma\).
This is \cite[Theorem~3.2.10(ii)]{Thu05}; properness removes the
boundary exception in that theorem.  It asserts the identity for real
piecewise affine \(F\), not just rational model functions.  The
composite \(F\circ\tau\) is lisse in Thuillier's sense.

If \(\mathfrak L\) determines a rational formal reference metric,
its curvature is \(\sum_v d_v\delta_v\) with the degrees in the
statement.  For CLD this follows from
\cite[Theorem~18.5.5 and Proposition~18.5.6(a)--(c)]{CLD25}:
the generic fibre is smooth and the strictly semistable special fibre
is geometrically reduced, so the multiplicity, defect, and graded
residue factors there all equal one.  The local curve comparison
identifies this with Thuillier curvature; details of that comparison
and its field normalization are recalled in Step~5 below.
Consequently \(\sum_vd_v=d\), by the degree theorem
\cite[Theorem~18.3.11]{CLD25}.  A degree on an integral component
with constant field \(k_v\) is
\([k_v:\widetilde K]\) times its degree over \(k_v\).
It must not be divided by this factor in the curvature measure.

\emph{2. The continuous potential and its curvature.}
Apply \cref{lem:ch31-green-kernel} with
\(\nu=\mu-d\delta_q\).  Choose finite positive atomic measures
\(\mu_j\) on \(G\), of mass exactly \(d\), converging weakly to
\(\mu\).  For instance partition each edge into intervals of mesh
tending to zero, assign each interval's mass to one point in it, and
keep vertex masses at the vertices.  Only real masses are required.
The corresponding potentials \(U_j=U_{\mu_j-d\delta_q}\) converge
uniformly to \(U\).  Each \(U_j\) is real piecewise affine; its
breakpoints may be type \(2\) or type \(3\).

By \cref{prop:ch29-curve-point-reference} and \eqref{eq:ch29-curve-retraction-current},
the lisse metric
\(\lVert\cdot\rVert_q\mathrm e^{-U_j\circ\tau}\) has curvature
\(\mu_j\).  Its local weights are subharmonic by
\cite[Theorem~3.4.12(ii)]{Thu05}.  Their locally uniform limits are
subharmonic: on a compact analytic test domain, compare with a harmonic
majorant and use uniform convergence on the domain and its boundary.
For a compactly supported lisse test function \(a\), the measure
\(\mathrm d\mathrm d^ca\) is finite.  Thus uniform convergence
also implies convergence of the local curvature currents by
\[
 \left|\left\langle\mathrm d\mathrm d^c(h_j-h),a\right\rangle\right|
       \leq\lVert h_j-h\rVert_\infty
                         \lVert\mathrm d\mathrm d^ca\rVert_{\mathrm{TV}}.
\]
The curvature of \eqref{eq:ch29-weighted-curve-solution} is therefore
\(\mu\).  Lisse tests are dense among continuous compactly supported
tests \cite[Proposition~3.2.7]{Thu05}, so the equality is an equality
of Radon measures.  Uniqueness follows from the same global harmonic
maximum principle as in \cref{prop:ch29-curve-point-reference}.
The identical argument starting with \(\mu_0\) proves
\eqref{eq:ch29-curve-model-reference-solution}; nefness of that
reference is not needed, since the approximating target curvatures
\(\mu_j\) are positive.

\emph{3. Rational approximation upstairs only.}
Put \(F=\widehat{K^{\mathrm a}}\).  By
\cite[Corollary~3.4.5]{Thu05}, the pulled-back continuous metric has
Thuillier-subharmonic weights.  The curve remains smooth, proper and
connected because it is geometrically connected.  Apply \cref{thm:p063-current-psh-ordinary-regularization} over the
\emph{algebraically closed} field \(F\): these weights are locally
CLD-psh.  The positive branch of \cref{thm:p081-corrected-curve-comparison} gives semipositive rational
formal metrics on \(L_F\) converging uniformly to the pullback.
No continuous-curve comparison over the nonclosed field is assumed.

\emph{4. Descent of approximation, without a common semistable model.}
On a proper algebraic variety formal and algebraic semipositive model
metrics agree by \cite[Theorem~1.1]{GM19}; thus the actual formal
approximants below may be used as rational algebraic model competitors.
Choose any rational model metric \(\lVert\cdot\rVert_b\) on \(L\);
its existence is \cite[\S2.1 and Proposition~2.10]{GM19}.
It need not be nef and need not be determined on \(\mathfrak C\).
Write the constructed metric as
\(\lVert\cdot\rVert_b\mathrm e^{-u}\), and let
\(\theta=c_1(L,\lVert\cdot\rVert_b)\).
The approximation in Step~3 gives model functions \(\varphi_j\)
over \(F\), with
\(\lVert\varphi_j-u_F\rVert_\infty<1/j\), such that
\(\theta_F+\mathrm d\mathrm d^c\varphi_j\) is nef.

By \cite[Proposition~2.18(b)]{GM19}, each rational piecewise linear
function \(\varphi_j\) is defined over a finite subextension
\(E_j/K\) inside \(K^{\mathrm a}\).  The corresponding metric
is semipositive over \(E_j\): represent it formally, pull its model
to \(F\), use the model-independence of semipositivity
\cite[Proposition~3.5]{GM19}, and reflect nefness through the residue
field extension by \cite[Lemma~3.3]{GM19}.  Take a finite normal
extension \(N_j/K\) containing \(E_j\).  It can include a purely
inseparable part.  Choose
\[
 c_j\in\mathbf Q(-\log|E_j^\times|),\qquad 1/j<c_j<2/j.
\]
Such constants exist because this nonzero divisible subgroup of
\(\mathbf R\) is dense; they are rational model constants of zero
curvature.  After extending to \(N_j\), the function
\(\psi_j=\varphi_j-c_j\) satisfies
\[
       u_{N_j}-3/j<\psi_j\leq u_{N_j}.
\]
The bounds descend from \(F\) by surjectivity of the scalar projection.

Here use \emph{only} the rational-model envelope
\[
 P^{\mathcal D_{\mathbf Q}}_\theta(u)
 =\sup\{a\in\mathcal D_{\mathbf Q}(C):
                 \theta+\mathrm d\mathrm d^ca\text{ nef},\ a\leq u\}.
\]
The actual competitor \(\psi_j\) gives
\[
 u_{N_j}-3/j
 \leq P^{\mathcal D_{\mathbf Q}}_{\theta_{N_j}}(u_{N_j})
 \leq u_{N_j}.
\]
For the finite \emph{normal} extension, precisely
\cite[Proposition~2.11]{GJKM19} identifies the middle term with the
pullback of \(P^{\mathcal D_{\mathbf Q}}_\theta(u)\).  Surjectivity
and \(j\to\infty\) give
\(P^{\mathcal D_{\mathbf Q}}_\theta(u)=u\).
At each point a model competitor is within \(\varepsilon\) of the
continuous function \(u\).  Its strict comparison remains true in
a neighborhood.  Choose a finite subcover and take the maximum of
the corresponding finitely many competitors.  Finite maxima preserve
the model nef class \cite[\S2.5]{GJKM19}, so this produces one model
function \(a_\varepsilon\) with
\(u-\varepsilon<a_\varepsilon\leq u\) everywhere.
This is the elementary argument of
\cite[Proposition~2.9(8)]{GJKM19}; no envelope continuity theorem is
an input.  It proves the required rational formal approximation over
\(K\), without using \cite[Theorem~3.1 or Proposition~3.10]{GJKM19}, without
identifying a real-linear span with a rational-model class, and without
asserting that a Green potential with rational masses has rational slopes.

\emph{5. The measure dictionary and scalar normalization.}
For an algebraically closed ground field, the curvature of a rational
formal metric agrees with Thuillier curvature.  One may check this
on a common semistable refinement: at type-\(2\) points the local
lisse comparison is \cite[Proposition~4.2]{Wanner19}; on an open
annulus it is the one-dimensional distributional second derivative;
at type-\(1\) and type-\(4\) points a lisse weight is locally
constant.  Wanner's model-function version is
\cite[Theorem~3.28]{Wanner19}.  All these uses of Wanner are over
an algebraically closed field.  Clearing the rational denominator
and working in local frames gives the metric statement.

The local weights of a rational formal metric over the original
\(K\) are themselves Thuillier-lisse; this fact does not require the
existence of a semistable model over \(K\) determining that metric.
Indeed, clear the rational denominator, so that the metric is induced
by a formal line bundle extending \(L^{\otimes m}\).  On a formal
trivializing chart let \(\varepsilon\) be its frame, of norm one.
After restricting to a strict affinoid domain where \(L\) has an invertible
frame \(e\), write \(\varepsilon=a e^{\otimes m}\), with \(a\) an
invertible analytic function.  Then
\[
                   -\log\lVert e\rVert=\frac1m\log|a|.
\]
These formulas hold on a covering in the analytic \(G\)-topology.
By \cite[Proposition~3.2.5]{Thu05}, the sheaf of Thuillier-lisse
functions on that site is the sheaf associated to real linear
combinations of logarithms of analytic units.  The continuous local
weights above are therefore lisse.  This argument applies equally
to the rational formal reference and to the rational formal
approximants.

To descend this identity for a rational model metric over \(K\),
use ordinary measure pushforward, with no field-degree multiplier.
For Thuillier it is \cite[Theorem~3.2.10(i)]{Thu05}, applied to
the local lisse weights.  For CLD it is
\cite[Corollary~18.3.3(a)]{CLD25}.  The skeletal explanation of this
normalization is \cite[Proposition~2.2.21]{Thu05}: if \(p:G_E\to G\)
is scalar extension, its integral affine edge coordinates are locally
unchanged, and
\begin{equation}
\label{eq:ch29-curve-scalar-weights}
 c_e=\sum_{e'\text{ above }e}c_{e'},
 \qquad
 p_*\Delta_{c'}(F\circ p)=\Delta_cF.
\end{equation}
The first equality is the residue-degree identity for the finite
\emph{reduced} fibre of an \'etale node chart.  At a vertex the sum
runs over all lifted points and all lifted tangent directions.
It is a sum identity, not equality of each individual component degree.
This scalar projection must not be confused with a finite analytic
morphism between curves over the \emph{same} ground field, whose
push--pull formula includes its degree
\cite[\S3.2.3]{Thu05}, \cite[Corollary~18.3.3(b)]{CLD25}.

Apply the model identity to the semipositive approximants from Step~4.
Their positive curvature measures all have mass \(d\).  On local
lisse test functions, uniform convergence of the weights gives
convergence to the constructed Thuillier curvature \(\mu\), as in
Step~2.  A finite partition of unity and density of lisse tests,
with the common mass bound, promote this to weak convergence on
\(C^{\mathrm{an}}\).  On the CLD side the same uniform convergence
gives convergence of curvature currents by their distributional
definition; the uniform-limit Chambert--Loir measure is the same
limit.  Hence both curvature measures of the final semipositive
metric are \(\mu\).

Finally, model metrics remain model and nef after every complete
scalar extension \cite[Proposition~2.18(a) and Lemma~3.3]{GM19}.
Surjectivity of scalar projection preserves the uniform distance,
so the pulled-back approximants converge to the pulled-back metric.
This proves the last assertion.  It does not claim continuity of
potentials for arbitrary nonskeletal measures.
\end{proof}


\begin{remark}[Source, date, and the division of labour]
\label{rem:ch29-curve-source}
For an ample line bundle, the envelope assertion of
\cref{thm:curves-arbitrary-field-continuity} is
\cite[Theorem~1.1]{GJKM19}, published in 2019.
The proof of \cref{thm:ch29-real-class-curve-model-envelope} treats
the specified rational-model envelope for real closed model forms;
it keeps the numerical class, the rational competitor class, and the
existence of an actual competitor separate.
The proof of the equation assertion uses the weighted graph kernel,
an original-field lisse reference metric, and the positive-degree
comparison \cref{thm:p081-corrected-curve-comparison}, with
\cref{thm:p063-current-psh-ordinary-regularization} used only over the
completed algebraic closure.  Finite normal descent includes the
inseparable case.  Neither envelope continuity nor the assertion
\(P_\theta(u)=u\) is used to infer semipositivity without competitors.
The higher-dimensional failure of the retraction statement is the
Burgos Gil--Sombra appendix to \cite{GJKM19}.
No assertion about continuous potentials for arbitrary nonskeletal
measures follows here.
\end{remark}

\begin{remark}[Where nefness and ampleness are used]
\label{rem:ch29-curve-nef-ample}
Positive degree, equivalently ampleness on the curve, is retained in
both assertions of \cref{thm:curves-arbitrary-field-continuity}.
Matching the total masses is what makes
\(\mu-d\delta_q\), or \(\mu-\mu_0\) when a model reference is actually
given, have mass zero.  The lisse reference is not asserted to be rational
formal.  In degree zero \cref{thm:ch29-real-class-curve-model-envelope}
requires an actual rational model competitor; generic nefness alone
does not supply one.
\end{remark}
